\section{STM32 CubeIDE}
STM32CubeIDE jest typem zintegrowanego środowiska programistycznego utworzonego przez firmę STMicroelectronics przeznaczone do projektowania aplikacji. Posiada główną funkcję programowania i debugowania kodów ładowanych na mikrokontrolery z rodziny STM32. Jest to narzędzie oparte na platformie Eclipse integrujące w jednym pakiecie wszystkie potrzebne funkcje do tworzenia systemów wbudowanych, co ułatwia testowanie aplikacji. \newline
\noindent Jednym z kluczowych elementów STM32CubeIDE jest STM32CubeMX, czyli graficzne narzędzie służące do konfiguracji mikrokontrolera, które umożliwia łatwe ustawianie peryferiów w pliku $.ioc$, np. zegarów systemowych i przypisywanie pinów. Generuje kod inicjalizacyjny dzieląc go na wszystkie potrzebne pliki dbając o porządek w kodzie użytkownika oraz pozwala bezpośrednio edytować go w środowisku STM32CubeIDE przyspieszając tworzenie prototypów. Ponadto środowisko posiada zaawansowany edytor kodu obsługujący języki C i C++ z funkcjami takimi jak podpowiadanie składni i automatyczne uzupełnianie kodu. Dodatkowo, STM32CubeIDE umożliwia debugowanie aplikacji w czasie rzeczywistym za pomocą interfejsów SWD z całą gamą narzędzi śledzących wartości poszczególnych zmiennych. Debugowanie obejmuje też śledzenie wykonania kodu, analizę błędów, przeglądanie rejestrów oraz pamięci mikrokontrolera, a także ustawianie punktów kontrolnych nazywanych potocznie breakpointami. \newline
\noindent STM32CubeIDE wspiera także RTOS - Real-Time Operating System, co ułatwia projektowanie aplikacji działających w czasie rzeczywistym. Umożliwia to zarządzanie wieloma procesami pozwalając rozbudowywać funkcjonalność programów. Narzędzie integruje również analizator sygnałów oraz debugowanie komunikacji UART przekładając się na możliwość monitorowania przesyłanych danych oraz optymalizację aplikacji.

\newpage

